% !TeX root = ../sections/02_exercises.tex
\subsection{2-B-4.}
\begin{exercise}
	数列 \(\{a_n\}_{n\in\mathbb{N}}\) の部分列
	\[
	\{a_{2k-1}\}_{k\in\mathbb{N}}
	\quad\text{および}\quad
	\{a_{2k}\}_{k\in\mathbb{N}}
	\]
	が共に同一の値 \(\alpha\in\mathbb{R}\) に収束するとき，
	\(\{a_n\}_{n\in\mathbb{N}}\) は \(\alpha\) に収束することを証明せよ。
\end{exercise}

\begin{background}
	数列 \(\{a_n\}\) は，奇数番目の項と偶数番目の項に分けられる。
	すなわち，任意の自然数 \(n\) は
	\[
	n=2k-1
	\quad\text{または}\quad
	n=2k
	\]
	のどちらか一方で表される。
	
	したがって，奇数番目の部分列
	\[
	\{a_{2k-1}\}
	\]
	と偶数番目の部分列
	\[
	\{a_{2k}\}
	\]
	がともに同じ極限 \(\alpha\) に収束するならば，
	十分後ろの項は，奇数番目であっても偶数番目であっても
	\(\alpha\) に近い。
	
	これを \(\varepsilon\)-\(N\) 論法で書けば，
	元の数列全体の収束が示される。
\end{background}

\begin{strategy}
	任意に \(\varepsilon>0\) を取る。
	
	仮定より
	\[
	a_{2k-1}\to\alpha
	\]
	であるから，ある \(K_1\in\mathbb{N}\) が存在して，
	\[
	k\geq K_1
	\quad\Longrightarrow\quad
	|a_{2k-1}-\alpha|<\varepsilon
	\]
	となる。
	
	また，
	\[
	a_{2k}\to\alpha
	\]
	であるから，ある \(K_2\in\mathbb{N}\) が存在して，
	\[
	k\geq K_2
	\quad\Longrightarrow\quad
	|a_{2k}-\alpha|<\varepsilon
	\]
	となる。
	
	ここで
	\[
	N=\max\{2K_1-1,\ 2K_2\}
	\]
	とおく。
	\(n\geq N\) を任意に取る。
	
	もし \(n\) が奇数ならば，
	ある \(k\in\mathbb{N}\) により
	\[
	n=2k-1
	\]
	と書ける。
	さらに \(n\geq 2K_1-1\) なので \(k\geq K_1\) である。
	したがって
	\[
	|a_n-\alpha|
	=
	|a_{2k-1}-\alpha|
	<\varepsilon
	\]
	である。
	
	もし \(n\) が偶数ならば，
	ある \(k\in\mathbb{N}\) により
	\[
	n=2k
	\]
	と書ける。
	さらに \(n\geq 2K_2\) なので \(k\geq K_2\) である。
	したがって
	\[
	|a_n-\alpha|
	=
	|a_{2k}-\alpha|
	<\varepsilon
	\]
	である。
	
	よって，任意の \(n\geq N\) に対して
	\[
	|a_n-\alpha|<\varepsilon
	\]
	が成り立つ。
	したがって
	\[
	a_n\to\alpha
	\]
	である。
\end{strategy}